% NSE Data Preparation -- Method-of-Snapshots Compression
% Standalone document: compile with  pdflatex DATA_PREPARATION.tex
\documentclass[11pt]{article}

\usepackage[margin=1in]{geometry}
\usepackage{amsmath,amssymb,amsthm}
\usepackage{mathtools}
\usepackage{booktabs}
\usepackage{array}
\usepackage{enumitem}
\usepackage{xcolor}
\usepackage[colorlinks=true,linkcolor=blue,urlcolor=blue]{hyperref}

% --- notation shortcuts -----------------------------------------------------
\newcommand{\R}{\mathbb{R}}
\newcommand{\Xtr}{X_{\mathrm{tr}}}
\newcommand{\Xte}{X_{\mathrm{te}}}
\newcommand{\Vtr}{V_{\mathrm{tr}}}
\newcommand{\Vte}{V_{\mathrm{te}}}
\newcommand{\Atr}{A_{\mathrm{tr}}}
\newcommand{\Ate}{A_{\mathrm{te}}}
\newcommand{\sref}{s_{\mathrm{ref}}}
\newcommand{\Cred}{C_{\mathrm{red}}}
\newcommand{\yref}{y_{\mathrm{ref}}}
\newcommand{\floortr}{\mathrm{floor}_{\mathrm{tr}}}
\newcommand{\floorte}{\mathrm{floor}_{\mathrm{te}}}
\newcommand{\norm}[1]{\left\lVert#1\right\rVert}
\newcommand{\Fnorm}[1]{\left\lVert#1\right\rVert_F}
\newcommand{\code}[1]{\texttt{#1}}
\DeclareMathOperator{\rank}{rank}
\DeclareMathOperator{\spanop}{span}
\DeclareMathOperator{\diag}{diag}

\title{\bfseries NSE Data Preparation:\\[2pt]
\large Method-of-Snapshots Compression of the K\'arm\'an-Vortex State}
\author{greedyopm experiments}
\date{}

\begin{document}
\maketitle

\begin{abstract}
This note documents, with full mathematical detail, how the large transient
Navier--Stokes (K\'arm\'an vortex street, $\mathrm{Re}=30$) dataset is prepared
for the method-comparison experiment. The aim is to make a $d=41{,}682$-dimensional
state problem cheap enough to fit \emph{6 methods $\times$ 10 reduced dimensions}
repeatedly, without ever working in the full state dimension and without losing
accuracy. The implementation is \code{prep\_nse.py}; the consumer is
\code{nse\_compare.py}.
\end{abstract}

\section{Pipeline overview}
\begin{center}
\begin{tabular}{@{}llp{0.32\linewidth}@{}}
\toprule
step & file & output \\
\midrule
1. assemble raw snapshots & \code{compile\_nse.py} & \code{data/nse\_snapshots.npz} \\
2. \textbf{compress (this note)} & \code{prep\_nse.py} & \code{data/nse\_compressed.npz} \\
3. recover DOF coordinates & \code{recover\_dof\_coords.py} & \code{coords\_stored.npy} \\
4. compare methods & \code{nse\_compare.py} & \code{plots/} \\
\bottomrule
\end{tabular}
\end{center}

\section{The data and why the raw form is expensive}
\code{compile\_nse.py} assembles the state matrix and the sensor (output) matrix
\begin{equation}
V \in \R^{d\times n},\qquad C \in \R^{q\times d},\qquad y(t)=C\,v(t),
\end{equation}
with the concrete sizes
\begin{equation}
d = 41{,}682\ \text{interior velocity DOFs},\qquad
n = 2000\ \text{snapshots},\qquad
q = 6\ \text{sensor outputs}.
\end{equation}
Raw $V$ is stored as float32 ($\approx 333$\,MB). The comparison fits POD,
Greedy-QM, Geelen-QM, Greedy-Fixed-OPM, and two Greedy-OPM variants, each across
$r=1,\dots,10$. Several of these perform repeated SVD-like work whose cost scales
with the \emph{state dimension} $d$. Carrying a $41{,}682$-row matrix through all
of those fits is the bottleneck --- it is the linear algebra, not the physics,
that is expensive.

\section{Train/test split and centering}
Even columns become training, odd columns become test. The \emph{training} mean
is subtracted from both sets, so the test set never uses its own statistics:
\begin{equation}
\sref = \frac{1}{n_{\mathrm{tr}}}\sum_{j}(\Vtr)_{:,j},\qquad
\Xtr = \Vtr - \sref\mathbf 1^\top \in \R^{d\times 1000},\qquad
\Xte = \Vte - \sref\mathbf 1^\top,
\end{equation}
with $n_{\mathrm{tr}}=n_{\mathrm{te}}=1000$.

\section{Key observation: an orthonormal change of coordinates is free}
The training columns span a subspace of dimension at most $n_{\mathrm{tr}}=1000\ll d$.
Let $U\in\R^{d\times K}$ be an orthonormal basis of that span, with thin SVD
\begin{equation}
\Xtr = U\,\Sigma\,W^\top,\qquad
\Sigma=\diag(\sigma_1,\dots,\sigma_K),\quad
W\in\R^{1000\times K},\quad U^\top U = I_K .
\end{equation}
Every method compared here builds its reduced basis \emph{inside} $\spanop(U)$ and
reconstructs the state through linear-plus-nonlinear maps of the projected
coordinates. Because $U$ is orthonormal it is an \emph{isometry}: for vectors
$x,x'$ in the training span with coordinates $a=U^\top x$, $a'=U^\top x'$,
\begin{equation}
\langle x,x'\rangle = \langle Ua,Ua'\rangle = \langle a,a'\rangle,
\qquad \norm{x}=\norm{a}.
\end{equation}
Norms, inner products, and singular values are therefore \textbf{identical}
whether evaluated on the full vector $x\in\R^{d}$ or on its coordinate vector
$a\in\R^{K}$. Hence running any of these methods on the coordinate matrix
\begin{equation}
A = U^\top X \in \R^{K\times n}
\end{equation}
reproduces \emph{exactly} the same reconstruction errors and the same selected
modes (up to the fixed rotation $U$) as running on the full $X$ --- but in a
$K$-dimensional space rather than a $41{,}682$-dimensional one. For this dataset
the $\mathrm{Re}=30$ limit cycle is genuinely low rank:
\begin{equation}
\boxed{K = 408}\qquad\text{(after the energy truncation of \S\ref{sec:floors})},
\end{equation}
so the working dimension shrinks by roughly $100\times$.

\section{Method of snapshots: coordinates without ever forming \texorpdfstring{$U$}{U}}
The basis $U$ is a $d\times K$ object and is \textbf{never materialized}. Instead
we use the small Gram matrix
\begin{equation}
G = \Xtr^\top \Xtr \in \R^{1000\times1000},\qquad
G = \Xtr^\top\Xtr = W\Sigma(U^\top U)\Sigma W^\top = W\,\Sigma^2\,W^\top .
\end{equation}
A single $1000\times1000$ symmetric eigenproblem yields the right singular vectors
$W$ and the singular values $\sigma_i=\sqrt{\lambda_i}$. The coordinates then
follow algebraically.

\paragraph{Training coordinates (exact, since $U^\top U=I$):}
\begin{equation}
\boxed{\;\Atr = U^\top \Xtr = U^\top U\,\Sigma W^\top = \Sigma W^\top
      = \diag(\sigma)\,W^\top\;}
\end{equation}

\paragraph{Test coordinates:} using $U=\Xtr W\Sigma^{-1}$ from the SVD, so that
$U^\top=\Sigma^{-1}W^\top \Xtr^\top$,
\begin{equation}
\boxed{\;\Ate = U^\top \Xte = \Sigma^{-1} W^\top\bigl(\Xtr^\top \Xte\bigr)\;}
\end{equation}
The cross-Gram $\Xtr^\top \Xte \in \R^{1000\times1000}$ is the only other
$d$-touching product. In total, the \emph{only} operations that ever see the full
dimension $d$ are the two Gram products $\Xtr^\top\Xtr$ and $\Xtr^\top\Xte$, each
computed once. The resulting coordinate matrices are
\begin{equation}
\Atr,\Ate \in \R^{408\times1000}\qquad(\approx 6.5\ \text{MB total}).
\end{equation}

\section{Lossless representation and the two error floors}
\label{sec:floors}

\subsection{Training floor}
In exact arithmetic with no truncation, $K=\rank(\Xtr)=n_{\mathrm{tr}}$ and
$\spanop(U)$ contains every training column exactly, so $\Fnorm{\Atr}=\Fnorm{\Xtr}$
and the floor is zero. In practice we drop POD modes whose singular value falls
below $\sigma_1\sqrt{\texttt{ENERGY\_TOL}}$ with $\texttt{ENERGY\_TOL}=10^{-13}$,
which yields $K=408$ and a tiny residual tail:
\begin{equation}
\floortr = \sqrt{\Fnorm{\Xtr}^2 - \Fnorm{\Atr}^2}
         = \Bigl(\textstyle\sum_{i>K}\sigma_i^2\Bigr)^{1/2},
\qquad
\frac{\floortr}{\Fnorm{\Xtr}} = 7.8\times10^{-7}.
\end{equation}

\subsection{Test floor}
Test snapshots may have a component \emph{outside} the training span. By the
Pythagorean theorem (the projection $UU^\top\Xte$ is orthogonal to the residual):
\begin{equation}
\floorte = \sqrt{\Fnorm{\Xte}^2 - \Fnorm{\Ate}^2}
         = \Fnorm{\Xte - UU^\top \Xte},
\qquad
\frac{\floorte}{\Fnorm{\Xte}} = 8.0\times10^{-6}.
\end{equation}
This is the \textbf{irreducible} error that no training-span method can beat. It is
computed once and stored, so it is never mis-attributed to any individual method.
Both floors use $\Fnorm{A}^2=\Fnorm{U^\top X}^2$, the energy of $X$ captured inside
$\spanop(U)$; the difference from $\Fnorm{X}^2$ is the energy outside the kept basis.

\section{Folding the floor back into the reported full-state error}
A method returns a coordinate reconstruction $\hat A$, implicitly representing the
full-state reconstruction $\hat X = U\hat A$. Decompose the true centered state as
$X = UA + R$ with $R = X - UU^\top X \perp \spanop(U)$. Since
$U(A-\hat A)\in\spanop(U)$ is orthogonal to $R$,
\begin{equation}
\Fnorm{X - U\hat A}^2
  = \underbrace{\Fnorm{A-\hat A}^2}_{\text{evaluated in }K\text{ dim}}
  + \underbrace{\Fnorm{R}^2}_{=\ \mathrm{floor}^2}.
\end{equation}
This is exactly the reported relative full-state error:
\begin{equation}
\text{state\_rel}
  = \frac{\sqrt{\Fnorm{A-\hat A}^2 + \mathrm{floor}^2}}{\Fnorm{X}}
  = \frac{\Fnorm{X - U\hat A}}{\Fnorm{X}},
\end{equation}
with $\Fnorm{X}$ stored as \code{nrm\_train}/\code{nrm\_test}. The full
$d=41{,}682$-dimensional error is thus reported faithfully while every fit and
every reconstruction stay in $K=408$ dimensions.

\section{Sensor outputs --- also computed without \texorpdfstring{$U$}{U}}
Because $v=\sref + Ua$, the linear output map reduces to
\begin{equation}
y = Cv
  = \underbrace{C\sref}_{\yref\in\R^{6}}
  + \underbrace{C U}_{\Cred\in\R^{6\times K}}\,a,
\qquad
\Cred = CU = (C\Xtr)\,W\,\Sigma^{-1}.
\end{equation}
Again $U$ never appears explicitly: $C\Xtr\in\R^{6\times1000}$ is tiny. The
ground-truth outputs $Y=CV$ are stored directly
($Y_{\mathrm{tr}},Y_{\mathrm{te}}\in\R^{6\times1000}$), so the output relative error
\begin{equation}
\text{out\_rel}
  = \frac{\Fnorm{Y - (\yref\mathbf 1^\top + \Cred\hat A)}}{\Fnorm{Y}}
\end{equation}
needs \textbf{no floor correction}: $Y$ already carries the full-state
information, including the out-of-span component, since
$\yref + \Cred\hat A = C(\sref + U\hat A) = C\hat v$ and the difference
$Y-\hat y = C(X - U\hat A)$ includes $CR$ automatically.

\section{What the downstream code ever sees}
After preparation, \code{nse\_compare.py} loads only the compressed arrays:
\begin{center}
\begin{tabular}{@{}lll@{}}
\toprule
array & shape & role \\
\midrule
\code{A\_train}, \code{A\_test} & $408\times1000$ & POD coordinates (compressed state) \\
\code{C\_red}                   & $6\times408$    & reduced output map $CU$ \\
\code{y\_ref}                   & $6$             & output of the mean state $C\sref$ \\
\code{Y\_train}, \code{Y\_test} & $6\times1000$   & ground-truth outputs $CV$ \\
\code{sigma}                    & $408$           & singular values (diagnostics) \\
\code{floor\_train}, \code{floor\_test} & scalar & irreducible error floors \\
\code{nrm\_train}, \code{nrm\_test}     & scalar & $\Fnorm{X}$ for relative errors \\
\code{K}                        & scalar          & $408$ \\
\bottomrule
\end{tabular}
\end{center}
The $41{,}682$-dimensional state is \emph{absent from the inner loop entirely}.

\paragraph{Cost summary.}
\begin{center}
\begin{tabular}{@{}lcc@{}}
\toprule
 & raw & compressed \\
\midrule
working state dim & $41{,}682$ & $\mathbf{408}$ \\
training matrix & $41{,}682\times1000$ & $408\times1000$ \\
on-disk size & $\approx 333$\,MB & $\approx 6.5$\,MB \\
per-method linear-algebra cost & $\mathcal O(d\,n^2)$ & $\mathcal O(K\,n^2)$, $\approx 100\times$ less \\
\bottomrule
\end{tabular}
\end{center}

\section{Practical notes}
\begin{itemize}[leftmargin=1.4em]
\item \textbf{Centering} uses the \emph{training} mean for both sets, so the test
  split never peeks at its own statistics.
\item \textbf{\texttt{ENERGY\_TOL} $=10^{-13}$} is what turns the idealized
  ``$\floortr=0$'' into the actual $\sim 8\times10^{-7}$ --- negligible, but
  nonzero; lowering it keeps more modes (larger $K$, smaller floor), raising it
  keeps fewer.
\item \textbf{Reproducing the numbers:} the concrete values
  ($K=408$, $\floortr/\Fnorm{\Xtr}=7.8\times10^{-7}$,
  $\floorte/\Fnorm{\Xte}=8.0\times10^{-6}$, $\Fnorm{\Xtr}=898.59$,
  $\Fnorm{\Xte}=898.48$) come from the saved \code{data/nse\_compressed.npz}.
\end{itemize}

\end{document}
